\section*{Bab \ref{ch:b1:structures} --- Struktur Aljabar}

\begin{solution}{exo:b1:structures:1}
\emph{Ketertutupannya:} $x * y = 1 \iff x + y - xy = 1 \iff (1-x)(1-y) = 0$,
yang mustahil untuk $x, y \neq 1$. Memang kesamaan kuncinya adalah
\[
1 - x * y = (1 - x)(1 - y):
\]
\omterm{def:b1:logic:map}{pemetaan} $\varphi(x) = 1 - x$ mengirim $(E, *)$ ke $(\R^*, \times)$ dengan
$\varphi(x * y) = \varphi(x)\varphi(y)$ --- yaitu \omterm{def:b1:structures:morphism}{morfisma} yang \omterm{def:b1:logic:inj}{bijektif}. Semua
aksiomanya kini terangkut: keasosiatifan dan kekomutatifannya menyusul dari
keasosiatifan dan kekomutatifan $\times$; identitasnya adalah $\varphi^{-1}(1) = 0$ (periksa: $x *
0 = x$); dan invers $x$ adalah $\varphi^{-1}\bigl((1-x)^{-1}\bigr) =
1 - \frac{1}{1-x} = \frac{x}{x - 1}$ (yang bernilai $\neq 1$). Jadi $(E, *)$
adalah \omterm{def:b1:structures:group}{grup abelian}.
\end{solution}

\begin{solution}{exo:b1:structures:2}
\begin{enumerate}
  \item Ya: hasil kali bilangan positif tetap positif, identitasnya $1$, inversnya
        $\frac 1x$, dan keasosiatifannya diwarisi dari $\R^*$.
  \item Bukan: tidak tertutup ($1 + 1 = 2 \notin \{-1,0,1\}$).
  \item Ya: yaitu contoh yang baku.
  \item Bukan: tidak tertutup (ganjil $+$ ganjil $=$ genap), dan tanpa identitas (karena $0$
        genap).
\end{enumerate}
\end{solution}

\begin{solution}{exo:b1:structures:3}
Tulis $\mathrm{id}$, transposisi $\tau_{12}, \tau_{13},
\tau_{23}$ (yang menukarkan kedua titik yang dinamai), serta siklus $c =
(1\,2\,3)$ (yakni $1 \mapsto 2 \mapsto 3 \mapsto 1$) dan $c^2 =
(1\,3\,2)$. Tabel $\sigma\rho$ (dengan baris $\sigma$, kolom $\rho$,
dan $\rho$ diterapkan lebih dulu):
\begin{center}
\renewcommand{\arraystretch}{1.2}
\begin{tabular}{c|cccccc}
$\sigma\backslash\rho$ & $\mathrm{id}$ & $c$ & $c^2$ & $\tau_{12}$ &
$\tau_{13}$ & $\tau_{23}$ \\
\hline
$\mathrm{id}$ & $\mathrm{id}$ & $c$ & $c^2$ & $\tau_{12}$ &
$\tau_{13}$ & $\tau_{23}$ \\
$c$ & $c$ & $c^2$ & $\mathrm{id}$ & $\tau_{13}$ & $\tau_{23}$ &
$\tau_{12}$ \\
$c^2$ & $c^2$ & $\mathrm{id}$ & $c$ & $\tau_{23}$ & $\tau_{12}$ &
$\tau_{13}$ \\
$\tau_{12}$ & $\tau_{12}$ & $\tau_{23}$ & $\tau_{13}$ & $\mathrm{id}$
& $c^2$ & $c$ \\
$\tau_{13}$ & $\tau_{13}$ & $\tau_{12}$ & $\tau_{23}$ & $c$ &
$\mathrm{id}$ & $c^2$ \\
$\tau_{23}$ & $\tau_{23}$ & $\tau_{13}$ & $\tau_{12}$ & $c^2$ & $c$ &
$\mathrm{id}$ \\
\end{tabular}
\end{center}
Pasangan yang tak berkomutasi: $\tau_{12}\tau_{13} = c^2$ sedangkan
$\tau_{13}\tau_{12} = c$. (Untuk memeriksa satu entri:
$\tau_{12}\tau_{13}$ mengirim $1 \xmapsto{\tau_{13}} 3
\xmapsto{\tau_{12}} 3$, $3 \mapsto 1 \mapsto 2$, $2 \mapsto 2 \mapsto
1$: yaitu $1 \mapsto 3 \mapsto 2 \mapsto 1$, siklus $c^2 =
(1\,3\,2)$.)
\end{solution}

\begin{solution}{exo:b1:structures:4}
$H$: di sini $1 \in H$; dan untuk $z, w \in H$, $\abs{zw^{-1}} = \abs z / \abs w =
1$: jadi kriterianya berlaku. Untuk $\R_+^*$: serupa itu, dengan $\abs{xy^{-1}}$ digantikan oleh
kepositifannya. Gabungannya: $\iu \in H$ dan $2 \in \R_+^*$, tetapi $2\iu$ bermodulus
$2 \neq 1$ dan bukan bilangan real positif: jadi $2\iu \notin H \cup
\R_+^*$, sehingga gabungannya tidak tertutup --- bukan \omterm{def:b1:structures:subgroup}{subgrup} (sebagaimana diramalkan
\cref{exo:b1:structures:6}, karena tak satu pun \omterm{def:b1:structures:subgroup}{subgrupnya} memuat yang lain).
\end{solution}

\begin{solution}{exo:b1:structures:5}
\omterm{def:b1:structures:morphism}{Morfisma}: $\eu^{\iu(\theta + \varphi)} =
\eu^{\iu\theta}\eu^{\iu\varphi}$ (\cref{thm:b1:complex:funceq}).
\omterm{def:b1:structures:morphism}{Kernelnya}: $\eu^{\iu\theta} = 1 \iff \theta \in 2\pi\Z$, jadi $\ker f =
2\pi\Z \neq \{0\}$: sehingga tidak \omterm{def:b1:logic:inj}{injektif}. Petanya: setiap bilangan kompleks bermodulus satu berupa
$\eu^{\iu\theta}$ untuk suatu $\theta$ (dari bentuk kutubnya), jadi
$\operatorname{im} f = \mathbb{U}$, yaitu lingkaran satuan. Pembatasan
$f$ pada selang setengah terbuka yang panjangnya $2\pi$, seperti
$\intco{0}{2\pi}$ atau $\intoc{-\pi}{\pi}$, bersifat \omterm{def:b1:logic:inj}{injektif} (karena dua sudut
yang petanya sama berselisih kelipatan $2\pi$, dan hanya satu
wakil tiap kelasnya yang muat di selang itu); tak ada selang yang
lebih panjang yang berhasil, karena ia memuat dua titik berjarak
$2\pi$.
\end{solution}

\begin{solution}{exo:b1:structures:6}
Irisannya: $e \in H \cap K$, dan $x, y \in H \cap K$ memberikan
$xy^{-1}$ di $H$ maupun di $K$. Gabungannya: jika $H \subseteq K$ maka gabungannya
adalah $K$, yaitu \omterm{def:b1:structures:subgroup}{subgrup} (dan simetris dengan itu). Sebaliknya, andaikan tak satu pun
inklusinya berlaku: pilih $h \in H \setminus K$ dan $k \in K \setminus
H$, lalu andaikan $H \cup K$ sebuah \omterm{def:b1:structures:subgroup}{subgrup}; maka $hk \in H \cup K$.
Jika $hk \in H$, maka $k = h^{-1}(hk) \in H$: kontradiksi. Jika $hk \in
K$, maka $h = (hk)k^{-1} \in K$: kontradiksi. Jadi $H \cup K$ bukan
\omterm{def:b1:structures:subgroup}{subgrup}.
\end{solution}

\begin{solution}{exo:b1:structures:7}
Perhatikan dulu bahwa $x^2 = e$ berarti $x^{-1} = x$ untuk setiap $x$. Lalu untuk
$x, y \in G$:
\[
xy = (xy)^{-1} = y^{-1} x^{-1} = yx ,
\]
dengan memakai \cref{prop:b1:structures:rules}~(2). Jadi $G$ \omterm{def:b1:structures:group}{abelian}.
\end{solution}

\begin{solution}{exo:b1:structures:8}
Unit $\Z/18\Z$: yaitu kelas yang \omterm{cor:b1:arith:bezout}{saling prima} dengan $18 = 2 \times 3^2$:
$\overline 1, \overline 5, \overline 7, \overline{11}, \overline{13},
\overline{17}$. Invers $\overline 5$: karena $5 \times 11 = 55 = 3\times
18 + 1$, maka $\overline 5^{-1} = \overline{11}$.

$\overline 5 x = \overline 7$: kalikan dengan $\overline{11}$: $x =
\overline{77} = \overline 5$ (karena $77 = 4\times 18 + 5$). Penyelesaiannya
tunggal.

$\overline 6 x = \overline 3$: persamaan $6x \equiv 3 \pmod{18}$
berarti $18 \mid 6x - 3$. Tetapi $6x - 3 = 3(2x - 1)$ ganjil, sedangkan $18$
genap: dan bilangan genap tak dapat \omterm{def:b1:arith:divides}{membagi} bilangan ganjil. Jadi tak ada penyelesaian.

$\overline 6 x = \overline{12}$: $6x \equiv 12 \pmod{18} \iff x
\equiv 2 \pmod 3$: penyelesaiannya $x \in \{\overline 2, \overline 5,
\overline 8, \overline{11}, \overline{14}, \overline{17}\}$ --- enam
buah.
\end{solution}

\begin{solution}{exo:b1:structures:9}
\omterm{def:b1:logic:sets}{Himpunan} $\Z[\sqrt 2]$ memuat $0$ dan $1$, serta tertutup terhadap pengurangan
dan perkalian:
\[
(a + b\sqrt 2)(c + d\sqrt 2) = (ac + 2bd) + (ad + bc)\sqrt 2 ,
\]
jadi ia subring $\R$ (kekomutatifan, keasosiatifan, dan
sifat distributifnya diwarisi). Unitnya: $(1 + \sqrt 2)(-1 + \sqrt 2) =
2 - 1 = 1$, jadi $1 + \sqrt 2$ terbalikkan dengan invers $\sqrt 2 - 1
\in \Z[\sqrt 2]$. Pangkatnya $(1 + \sqrt 2)^n$ naik tegas
(karena bilangan pokoknya $> 1$), sehingga berbeda dua-dua, dan masing-masing
merupakan unit (karena $\bigl((1+\sqrt2)^n\bigr)^{-1} = (\sqrt 2 - 1)^n$): jadi \omterm{def:b1:structures:group}{grup}
unitnya tak hingga.
\end{solution}

\begin{solution}{exo:b1:structures:10}
$x + x = (x + x)^2 = x^2 + x^2 + x^2 + x^2 = 4x^2 = 4x$ --- sehingga $2x =
4x$, yang memberikan $2x = 0$, yakni $x + x = 0$ (jadi setiap unsurnya menjadi
inversnya sendiri terhadap penjumlahan). Lalu
\[
x + y = (x+y)^2 = x^2 + xy + yx + y^2 = x + xy + yx + y ,
\]
sehingga $xy + yx = 0$, yakni $xy = -yx = yx$ (dengan memakai $-z = z$). Jadi $A$
komutatif.

Contohnya: pada $\mathcal{P}(E)$, definisikan $A + B = (A \cup B) \setminus (A
\cap B)$ (yaitu selisih simetrisnya) dan $A \times B = A \cap B$. Kita
periksa: $(\mathcal{P}(E), +)$ adalah \omterm{def:b1:structures:group}{grup abelian} dengan identitas
$\emptyset$ dan setiap \omterm{def:b1:logic:sets}{himpunan} menjadi inversnya sendiri; $\cap$ bersifat asosiatif,
komutatif, dengan identitas $E$; dan sifat distributifnya $A \cap (B + C) = (A
\cap B) + (A \cap C)$ berlaku (karena sebuah unsur berada di ruas kiri jika dan hanya jika ia
berada di $A$ dan tepat di salah satu dari $B, C$). Lalu $A \cap A = A$: jadi setiap
unsurnya idempoten, sesuai tuntutannya.
\end{solution}

\begin{solution}{exo:b1:structures:11}
Tuliskan hipotesisnya untuk $n$ dan $n+1$:
\[
(xy)^{n+1} = x^{n+1} y^{n+1}
\quad\text{dan}\quad
(xy)^{n+1} = (xy)(xy)^n = xy\,x^n y^n .
\]
Dengan menyamakannya: $x^{n+1} y^{n+1} = x\,y\,x^n\,y^n$; coret $x$ di kiri
dan $y^n$ di kanan: maka $x^n y = y x^n$. Perhitungan yang sama satu
derajat lebih tinggi ($n+1$ dan $n+2$) memberikan $x^{n+1} y = y x^{n+1}$. Lalu
\[
y\,x^{n+1} = x^{n+1} y = x\,(x^n y) = x\,y\,x^n ,
\]
dan dengan mencoret $x^n$ di sebelah kanan pada $y x \cdot x^n = x y \cdot x^n$:
diperoleh $yx = xy$. Jadi $G$ \omterm{def:b1:structures:group}{abelian}.
\end{solution}

\begin{solution}{exo:b1:structures:12}
\begin{enumerate}
  \item Misalkan $f \colon \Z \to \Z$ aditif dan $a = f(1)$. Dengan
        induksi, $f(k) = ka$ untuk $k \in \N$, dan $f(-k) = -f(k) =
        -ka$: jadi $f$ adalah perkalian dengan $a$. Sebaliknya
        setiap \omterm{def:b1:logic:map}{pemetaan} $k \mapsto ak$ merupakan \omterm{def:b1:structures:morphism}{morfisma}: jadi \omterm{def:b1:structures:morphism}{morfisma}
        $(\Z,+) \to (\Z,+)$ tepat berupa perkalian dengan sebuah
        bilangan bulat tetap.
  \item Misalkan $f \colon \Q \to \Z$ sebuah \omterm{def:b1:structures:morphism}{morfisma}, $x \in \Q$ dan
        $n \in \N^*$. Maka
        \[
        f(x) = f\Bigl(\underbrace{\tfrac xn + \dots +
        \tfrac xn}_{n}\Bigr) = n\,f\Bigl(\frac xn\Bigr) ,
        \]
        sehingga bilangan bulat $f(x)$ habis dibagi setiap $n \geq 1$.
        Satu-satunya bilangan bulat semacam itu adalah $0$: jadi $f \equiv 0$.
\end{enumerate}
\end{solution}

\begin{solution}{pb:b1:structures:1}
\textbf{1.} Sebuah \omterm{def:b1:counting:objects}{permutasi} adalah bijeksi $\intint1n$, yakni
susunan-$n$ atas $n$ objek: dan ada $n!$ di antaranya
(\cref{thm:b1:counting:counts}). Dengan $\sigma = [2,3,1]$, $\tau =
[1,3,2]$: $\sigma\tau$ mengirim $1 \mapsto 1 \mapsto 2$, $2 \mapsto 3
\mapsto 1$, $3 \mapsto 2 \mapsto 3$: jadi $\sigma\tau = [2,1,3]$; sedangkan
$\tau\sigma$ mengirim $1 \mapsto 2 \mapsto 3$, $2 \mapsto 3 \mapsto
2$, $3 \mapsto 1 \mapsto 1$: jadi $\tau\sigma = [3,2,1] \neq
\sigma\tau$.

\textbf{2.} Misalkan $\gamma, \gamma'$ bertumpuan saling lepas $S, S'$.
Untuk $x \in S$: $\gamma'(x) = x$ dan $\gamma(x) \in S$, jadi
$\gamma\gamma'(x) = \gamma(x) = \gamma'\gamma(x)$. Simetris dengan itu
untuk $x \in S'$; dan kedua ruasnya menetapkan setiap $x \notin S \cup S'$. Jadi
$\gamma\gamma' = \gamma'\gamma$.

\textbf{3.} Untuk $i \in \intint1n$, nilai $i, \sigma(i),
\sigma^2(i), \dots$ hidup di dalam \omterm{def:b1:counting:card}{himpunan hingga}, jadi $\sigma^a(i) =
\sigma^b(i)$ untuk suatu $a < b$; lalu keinjektifannya memberikan $\sigma^{b-a}(i)
= i$: sehingga barisannya kembali ke $i$. Sebutlah \emph{orbit} $i$ sebagai
\omterm{def:b1:logic:sets}{himpunan} $\{i, \sigma(i), \dots, \sigma^{k-1}(i)\}$ dengan $k \geq 1$
yang minimal sedemikian sehingga $\sigma^k(i) = i$. Dua orbit yang bertemu pada satu
titik berimpit (karena keduanya adalah peta $\sigma$ ke depan dari titik
itu), jadi orbitnya mempartisi $\intint1n$; dan $\sigma$ bekerja pada setiap
orbit berukuran $k \geq 2$ sebagai siklus-$k$ $(i\ \sigma(i)\ \cdots\
\sigma^{k-1}(i))$ serta menetapkan singletonnya. Hasil kali siklus
yang saling lepas itu sama dengan $\sigma$ di mana-mana. Ketunggalannya: pada
sebarang penguraian menjadi siklus yang saling lepas, siklus yang melewati $i$ pastilah
$(i\ \sigma(i)\ \cdots)$ --- jadi siklusnya terpaksa menjadi
orbitnya beserta aksi yang terinduksi padanya.

\textbf{4.} Dengan mengikuti orbitnya: $1 \to 4 \to 2 \to 1$, $3 \to 5
\to 3$, $6 \to 7 \to 8 \to 6$:
\[
\sigma = (1\ 4\ 2)(3\ 5)(6\ 7\ 8) .
\]
Jika $\sigma = \gamma_1\cdots\gamma_r$ dengan siklus saling lepas berpanjang
$k_1, \dots, k_r$, maka kekomutasiannya (pertanyaan 2) memberikan
$\sigma^m = \gamma_1^m\cdots\gamma_r^m$, dan karena tumpuannya
saling lepas, $\sigma^m = \mathrm{id}$ jika dan hanya jika setiap $\gamma_i^m =
\mathrm{id}$, yakni jika dan hanya jika $k_i \mid m$ untuk setiap $i$ (karena siklus-$k$ berorde
$k$: $\gamma^m$ mengirim $a_1$ ke $a_{1 + (m \bmod k)}$). Nilai $m$
terkecil semacam itu adalah $\operatorname{lcm}(k_1, \dots, k_r)$. Di sini:
$\operatorname{lcm}(3, 2, 3) = 6$.

\textbf{5.} Terapkan ruas kanannya pada setiap titik, dimulai dari faktor
paling kanan. Di sini $a_1 \mapsto a_2$ oleh $(a_1\ a_2)$, lalu setiap faktor
berikutnya menetapkan $a_2$: jadi hasil bersihnya $a_1 \mapsto a_2$. Untuk $2 \leq i < k$:
$a_i$ tak tersentuh sampai $(a_1\ a_i)$ mengirimnya ke $a_1$, dan
faktor tepat berikutnya $(a_1\ a_{i+1})$ mengirim $a_1$ ke $a_{i+1}$, setelah
itu tak ada lagi yang menggerakkannya: jadi hasil bersihnya $a_i \mapsto a_{i+1}$. Akhirnya $a_k$
ditetapkan oleh semua faktornya kecuali yang paling kiri, yang mengirimnya ke $a_1$.
Inilah persis siklusnya. Karena setiap \omterm{def:b1:counting:objects}{permutasi} merupakan hasil kali
siklus (pertanyaan 3), ia hasil kali transposisi. Untuk
$\sigma$ pada pertanyaan 4:
\[
\sigma = (1\ 2)(1\ 4)\;(3\ 5)\;(6\ 8)(6\ 7),
\]
yaitu lima transposisi.

\textbf{6.} Induksi pada $b - a$. Untuk $b = a + 1$ kesamaannya
sepele (dengan $1 = 2\cdot1 - 1$ faktor). Untuk $b > a + 1$, periksa
langsung bahwa $(a\ b) = (a\ \ a{+}1)\,(a{+}1\ \ b)\,(a\ \ a{+}1)$:
ruas kanannya mengirim $a \mapsto a{+}1 \mapsto b \mapsto b$, $b
\mapsto b \mapsto a{+}1 \mapsto a$, $a{+}1 \mapsto a \mapsto a
\mapsto a{+}1$, dan menetapkan selebihnya. Menurut induksi, $(a{+}1\ \ b)$ adalah
hasil kali palindromik berisi $2(b - a - 1) - 1$ transposisi
bersebelahan, sehingga $(a\ b)$ berisi $2(b - a) - 1$ transposisi: yaitu bilangan
ganjil.

\textbf{7.} $N(\mathrm{id}) = 0$, $\varepsilon = +1$. Untuk $(i\ \
i{+}1)$, satu-satunya pasangan yang terbalik adalah $(i, i+1)$: jadi $N = 1$,
$\varepsilon = -1$. Untuk $[2, 3, 1]$: pasangan yang terbalik adalah $(1,
3)$ (dengan nilai $2 > 1$) dan $(2, 3)$ (dengan nilai $3 > 1$): jadi $N = 2$,
$\varepsilon = +1$.

\textbf{8.} Daftar nilai $\sigma$ dan $\sigma\tau$ hanya berbeda
oleh penukaran posisi $i$ dan $i + 1$. Untuk pasangan
posisi yang tak melibatkan $i, i+1$, tak ada yang berubah. Untuk $k < i$,
kedua pasangan $(k, i)$ dan $(k, i+1)$ saling menukarkan status
inversinya (karena kedua nilai yang sama dibandingkan dengan $\sigma(k)$, dengan
urutan posisi yang berlainan): sehingga sumbangan totalnya
tak berubah; demikian pula untuk $k > i + 1$. Satu pasangan yang tersisa, $(i,
i+1)$, membalik statusnya. Jadi $N(\sigma\tau) = N(\sigma) \pm 1$.

\textbf{9.} Misalkan $\tau = (a\ b)$ sebarang transposisi: menurut
pertanyaan 6 ia hasil kali sejumlah ganjil transposisi
bersebelahan, jadi mengalikannya di kanan dengan $\tau$ mengubah $N$ sebanyak
total yang ganjil (pertanyaan 8, yang diterapkan berulang): sehingga $\varepsilon(\sigma
\tau) = -\varepsilon(\sigma)$. Sekarang jika $\sigma = \tau_1\cdots
\tau_p$ (dengan transposisi), bangunlah ia dari identitasnya lewat $p$ kali
perkalian di kanan: maka $\varepsilon(\sigma) =
(-1)^p\varepsilon(\mathrm{id}) = (-1)^p$. Karena
$\varepsilon(\sigma)$ didefinisikan lewat inversinya --- tanpa bergantung
pada pemfaktoran apa pun --- maka paritas $p$ merupakan invarian
$\sigma$. \omterm{def:b1:structures:morphism}{Morfismanya}: dengan menulis $\sigma$ memakai $p$ dan $\sigma'$ memakai
$q$ transposisi, $\sigma\sigma'$ memakai $p + q$ transposisi:
sehingga $\varepsilon(\sigma\sigma') = (-1)^{p+q} =
\varepsilon(\sigma)\varepsilon(\sigma')$.

\textbf{10.} Siklus-$k$ adalah hasil kali $k - 1$ transposisi
(pertanyaan 5): jadi $\varepsilon = (-1)^{k-1}$. Untuk $\sigma$ yang umum
dengan orbit berukuran $k_1, \dots, k_r$ ($k_i \geq 2$) ditambah $f$
titik tetap, $c(\sigma) = r + f$ dan $n = k_1 + \dots + k_r + f$,
sehingga
\[
\varepsilon(\sigma) = \prod_{i=1}^r (-1)^{k_i - 1}
= (-1)^{\sum_i k_i - r} = (-1)^{n - f - r} = (-1)^{n -
c(\sigma)} .
\]

\textbf{11.} Di sini $\mathfrak A_n = \ker\varepsilon$ merupakan \omterm{def:b1:structures:subgroup}{subgrup} sebagai
\omterm{def:b1:structures:morphism}{kernel} sebuah \omterm{def:b1:structures:morphism}{morfisma} (\cref{def:b1:structures:morphism}). Tetapkan
sebuah transposisi $\tau_0$ (yang ada untuk $n \geq 2$). \omterm{def:b1:logic:map}{Pemetaan} $\sigma
\mapsto \sigma\tau_0$ adalah bijeksi $\mathfrak S_n$ (yang menjadi inversnya
sendiri) yang menukarkan $\mathfrak A_n$ dengan \omterm{def:b1:logic:sets}{himpunan} \omterm{def:b1:counting:objects}{permutasi}
ganjil (pertanyaan 9). Kedua \omterm{def:b1:logic:sets}{himpunan} itu mempartisi $\mathfrak S_n$
dan berukuran sama: jadi $\abs{\mathfrak A_n} = \frac{n!}2$.

\textbf{12.} \emph{Inversi} $[4, 1, 5, 2, 3, 7, 8, 6]$: dari
nilai $4$: atas $1, 2, 3$: tiga buah; dari $5$: atas $2, 3$: dua buah;
dari $7$: atas $6$: satu buah; dari $8$: atas $6$: satu buah. Jadi $N = 7$,
$\varepsilon = -1$. \emph{Tipe siklusnya}: $c = 3$ orbit, $n = 8$:
jadi $\varepsilon = (-1)^{8-3} = -1$. \emph{Cacah transposisinya}: lima
transposisi pada pertanyaan 5: jadi $(-1)^5 = -1$. Ketiganya sepakat.

\textbf{13.} $(a\ b)(a\ c)$ (dimulai dari paling kanan): $a \mapsto c
\mapsto c$; $c \mapsto a \mapsto b$; $b \mapsto b \mapsto a$:
yaitu siklus-$3$ $(a\ c\ b)$. Lalu $(a\ c\ b)(a\ c\ d)$: $a \mapsto c
\mapsto b$; $b \mapsto b \mapsto a$; $c \mapsto d \mapsto d$; $d
\mapsto a \mapsto c$: yaitu $(a\ b)(c\ d)$, sesuai klaimnya.

\textbf{14.} Misalkan $\sigma \in \mathfrak A_n$: menurut pertanyaan 9,
$\sigma = \tau_1\cdots\tau_{2m}$ dengan sejumlah genap
transposisi. Kelompokkan berpasangan secara berurutan
$\tau_{2i-1}\tau_{2i}$: jika keduanya sama, pasangannya menjadi
identitas dan lenyap; jika keduanya berbagi tepat satu titik, maka
kesamaan pertama pada pertanyaan 13 menuliskannya sebagai satu siklus-$3$;
dan jika keduanya saling lepas, kesamaan keduanya menuliskannya sebagai dua
siklus-$3$. Jadi $\sigma$ merupakan hasil kali siklus-$3$ (atau
identitasnya, yaitu hasil kali kosong --- dan untuk $n \geq 3$ juga $(1\ 2\
3)^3$).

\textbf{15.} $(1\ 2)(3\ 4) = (1\ 3\ 2)(1\ 3\ 4)$ (pertanyaan 13
dengan $a{=}1, b{=}2, c{=}3, d{=}4$). Untuk siklus-$5$: menurut
pertanyaan 5, $(1\ 2\ 3\ 4\ 5) = (1\ 5)(1\ 4)(1\ 3)(1\ 2)$, dan dengan
memasangkannya: $(1\ 5)(1\ 4) = (1\ 4\ 5)$, $(1\ 3)(1\ 2) = (1\ 2\ 3)$:
\[
(1\ 2\ 3\ 4\ 5) = (1\ 4\ 5)(1\ 2\ 3) .
\]
(Periksa pada $3$: $(1\ 2\ 3)$ mengirim $3 \to 1$, lalu $(1\ 4\ 5)$
mengirim $1 \to 4$: jadi hasil bersihnya $3 \to 4$, dan itu benar.)

\textbf{16.} Terapkan kedua ruasnya pada sebuah titik sembarang. Untuk $i =
\sigma(a_j)$: ruas kirinya memberikan $\sigma\bigl((a_1\ \dots\
a_k)(a_j)\bigr) = \sigma(a_{j+1})$ (dengan indeks modulo $k$), yang persis
dilakukan ruas kanannya pada $\sigma(a_j)$. Untuk $i$ yang tak berbentuk
demikian: $\sigma^{-1}(i)$ berada di luar tumpuannya, jadi ruas kirinya
menetapkan $i$, dan begitu pula ruas kanannya. Jadi keduanya sama di mana-mana.

\textbf{17.} Menggeser ubin di sel $c'$ ke dalam sel kosong
$c$ menukarkan isi sel $c$ dan $c'$ (karena ubin $9$, yaitu
sel kosongnya, pindah ke $c'$). Jika ubin $\sigma(i)$ duduk di sel $i$, maka
kedudukan barunya adalah $\sigma' = \sigma \circ (c\ c')$: isinya sama
kecuali bahwa sel $c, c'$ kini memuat isi yang sebelumnya di sel satunya. Menurut
pertanyaan 9, $\varepsilon(\sigma') = -\varepsilon(\sigma)$.

\textbf{18.} Sebuah langkah mengirim sel kosongnya ke sel yang bersebelahan: sehingga
baris atau kolomnya berubah tepat sebesar $1$, jadi jarak taksi
$d$ ke sel $9$ berubah sebesar $\pm1$, dan $(-1)^d$ membalik. Karena setiap
langkah membalik $\varepsilon(\sigma)$ sekaligus $(-1)^{d(\sigma)}$,
maka hasil kalinya $I(\sigma)$ tak berubah oleh setiap langkah: yaitu sebuah
invarian.

\textbf{19.} Kedudukan yang terpecahkan mempunyai $\varepsilon = +1$, $d = 0$:
jadi $I = +1$. Sasarannya (dengan ubin $7, 8$ tertukar dan sel kosong di tempat asalnya) adalah
transposisi isi sel $7$ dan $8$: sehingga $\varepsilon
= -1$, $d = 0$: jadi $I = -1$. Karena $I$ invarian dan kedua
nilainya berbeda, tak ada rangkaian langkah yang menghubungkan keduanya.

\textbf{20.} Kedudukan dengan sel kosong di tempat asalnya adalah \omterm{def:b1:counting:objects}{permutasi}
$8$ ubin di antara sel $1, \dots, 8$, yakni unsur
$\mathfrak S_8$; ia mempunyai $d = 0$, jadi $I = \varepsilon(\sigma)$.
Keterjangkauannya memaksa $I = +1$, yakni $\sigma \in \mathfrak A_8$; dan
konvers yang diterima tadi mengatakan seluruh $\mathfrak A_8$ tercapai. Cacahnya:
$\abs{\mathfrak A_8} = \frac{8!}2 = 20\,160$ (pertanyaan 11).

\textbf{21.} Menurut pertanyaan 20, susunan yang dapat dicapai dengan sel
kosong di tempat asalnya tepat membentuk $\mathfrak A_8$ --- khususnya sebuah
\emph{\omterm{def:b1:structures:subgroup}{subgrup}} $\mathfrak S_8$: jadi mengomposisikan dua pengacakan yang
terpecahkan, atau membalikkan salah satunya, tetap terpecahkan, dan itu jauh dari
kasatmata bila kita bernalar dengan teka-tekinya saja.

\textbf{22.} (a) Siklus-$3$ atas ubinnya dengan sel kosong di tempat asalnya:
$\varepsilon = +1$ (pertanyaan 10), $d = 0$, jadi $I = +1$: sehingga terjangkau
(menurut konvers yang diterima tadi) --- jadi kita dapat mendaurkan tiga ubin.
(b) Ubin $5$ dan sel kosongnya bertukar: kedudukannya adalah
transposisi isi sel $5$ dan $9$, sehingga
$\varepsilon = -1$; sedangkan sel kosongnya duduk di tengah, pada jarak taksi
$d = 2$ dari tempat asalnya, sehingga $(-1)^d = +1$ dan $I = -1$:
jadi tak terjangkau. Kita tak dapat begitu saja ``memarkir sel kosong di tengah''
sambil membiarkan ubin lainnya tersusun rapi.

\textbf{23.} Misalkan $f \colon \mathfrak S_n \to \{\pm1\}$ sebuah
\omterm{def:b1:structures:morphism}{morfisma}. Untuk sebarang dua transposisi $\tau, \tau'$, pertanyaan 16
memasok $\sigma$ dengan $\sigma\tau\sigma^{-1} = \tau'$ (yaitu petakan
kedua titik yang bergerak ke kedua titik yang lain; $n \geq 3$ menjamin ruang
untuk melakukannya, meskipun $n = 2$ pun sepele di sini). Maka $f(\tau') =
f(\sigma)f(\tau)f(\sigma)^{-1} = f(\tau)$ karena $\{\pm1\}$
\omterm{def:b1:structures:group}{abelian}: jadi $f$ konstan pada transposisi. Jika konstanta itu
$+1$, maka $f = 1$ pada semua hasil kali transposisi, yakni
di mana-mana (pertanyaan 5). Jika konstanta itu $-1$, maka $f(\sigma) =
(-1)^p = \varepsilon(\sigma)$ pada hasil kali $p$ transposisi.
Jadi $f \in \{1, \varepsilon\}$.

\textbf{24.} (i) Sifat \omterm{def:b1:structures:morphism}{morfisma} pada $\varepsilon$ dan
mesin \omterm{def:b1:structures:morphism}{kernelnya} memberi $\mathfrak A_n$ struktur \omterm{def:b1:structures:subgroup}{subgrup} beserta
ukurannya, dan penalaran bergaya \cref{prop:b1:structures:kernel}
mengalir di sepanjang pertanyaan 11 dan 21. (ii) Pencacahan:
$\abs{\mathfrak S_n} = n!$, argumen pemaruhan pada pertanyaan 11,
dan cacahan $20\,160$ pada pertanyaan 20 semuanya
\cref{ch:b1:counting} yang sedang bekerja. (iii) Pertanyaan 8--9 menyelesaikan
persoalan keterdefinisian yang sungguhan --- karena ``paritas banyaknya
transposisi'' mengandaikan bahwa paritas itu tidak bergantung
pada pemfaktorannya, persis seperti operasi $\Z/n\Z$ menuntut
ketakbergantungan pada wakilnya pada \cref{def:b1:structures:zn}.

\textbf{25.} Tanda adalah satu perhitungan bernilai $\{\pm1\}$,
yang sekali dibuktikan terdefinisi dengan baik, lalu mengerjakan tiga
tugas sekaligus: ke dalam, ia memotong $\mathfrak S_n$ menjadi dua dan
mengisolasi $\mathfrak A_n$ beserta pembangkit siklus-$3$-nya;
ke luar, ia memutuskan dalam satu baris sebuah pertanyaan (``dapatkah kedua
ubin ini ditukar?'') yang tak akan pernah dapat dituntaskan pencarian naif, karena
tak ada daftar hingga rangkaian langkah yang gagal yang membuktikan kemustahilannya; dan
secara struktural, ia adalah mesin tanda berselang-seling di dalam
rumus $\det A = \sum_\sigma \varepsilon(\sigma)\,
a_{1\sigma(1)}\cdots a_{n\sigma(n)}$ pada \cref{ch:b1:det}. Pola
itu --- yaitu carilah besaran yang terawetkan oleh setiap langkah dasar,
lalu hitung ia pada awal dan pada sasarannya --- adalah senjata baku
matematikawan melawan pertanyaan ``mungkinkah?'',
dan ia akan kembali setiap kali sebuah \omterm{def:b1:structures:group}{grup} bekerja pada sebuah \omterm{def:b1:logic:sets}{himpunan}
keadaan.

\end{solution}
